\section{Dataset inputs and provenance} 
\label{app:data} 
\setcounter{table}{0} 
Table~\ref{tab: dataset-inputs} summarises the inputs we assessed, drawn from the cache manifest of each dataset and the corresponding adapter scripts under \path{src/pace/datasets/adapters/}. The cache folder names (e.g.\ \texttt{IoT-23-v2}) should not be interpreted as the upstream version number. We provide the precise inputs and their hashes for reference in the supplementary file \path{revision-evidence/dataset-input-manifest.json}, where the source files and respective feature schemas are identified along with a hash check of the manifest. Hashes were independently checked for all 7 datasets. Note that the above does not attest to the validity of hashes for the underlying raw data. Where available, we additionally record content hashes of the campaigns, along with campaign input/preprocess hashes, which serve to distinguish our frozen campaign results. However, these again cannot be used to trace back to hashes of raw data.

{\small 
\begin{longtable}{>{\raggedright\arraybackslash}p{0.13\textwidth}>{\raggedright\arraybackslash}p{0.43\textwidth}>{\raggedright\arraybackslash}p{0.34\textwidth}} 
\caption{List of datasets considered, associated input files we consider in this work, and links/restrictions for the sources.}\label{tab:dataset-inputs}\\
\toprule Dataset & Evaluated input and adapter & Source and conditions\\\midrule\endfirsthead
\toprule Dataset & Evaluated input and adapter (continued) & Source and conditions\\\midrule\endhead
\bottomrule\endfoot 
IoT-23 & \texttt{iot\_23\_datasets\_small.tar.gz}, comprising preprocessed labelled Zeek traffic (as opposed to \textit{all} PCAP data). Version 1.0.0 listed in the dataset. See \texttt{iot23.py}. & \href{https: //www.stratosphereips.org/datasets-iot23}{Dataset website}; published on Zenodo under CC BY 4.0 \href{https://doi.org/10.5281/zenodo.4743746}{(record)}.\\
TON\_IoT & Data under \path{Processed\_datasets/Process\_Network\_dataset}, i.e.\ only the "network" CSV data and not the "telemetry" and OS-specific files. No version specified for cached input. See \texttt{ton\_iot.py}. & \href{https://research.unsw.edu.au/projects/toniot-datasets}{Dataset website}. Released for academic and non-profit research without consent of copyright holder. Commercial usage is conditional on written authorisation and source citation.\\
Gotham & Data uploaded to Zenodo in 2025; we consider 8 clients (each represented by a CSV under \path{extracted/processed/}). 7 classes are evaluated in the campaign (+1 class only for "client"). See \texttt{gotham.py}; see also Section~\ref{sec: experiment} regarding features calculated over entire seconds. & \href{https://doi.org/10.5281/zenodo.14502760}{Zenodo dataset} under CC BY 4.0. See also a recent 2026 dataset descriptor\,\cite{belarbi2026descriptor} which references the Zenodo release.\\
CTU-13 & \texttt{CTU-13-Dataset.tar.bz2}; we consider 12 of 13 scenarios (see above). See \texttt{ctu13.py}; labelled Argus data and exclude "Background" label. No other version information apparent. & \href{https: //www.stratosphereips.org/datasets-ctu13}{Dataset website}; \href{https://www.stratosphereips.org/datasets-faq}{FAQ} notes acknowledgement of the source but no details provided on redistribution of raw data.\\
N-BaIoT & As released in the original (2018) work, considering 9 devices. We consider the extracted benign and malware (Mirai & BASHLITE) data only. See \texttt{nbaiot.py}, noting the additional filtering applied of exact matching lines within files. & UCI ML repository entry \href{https: //doi.org/10.24432/C5RC8J}{(record)} under CC BY 4.0 \cite{nbaiotdata2018}; note this is different to the original dataset publication\,\cite{nbaiot2018}.\\
UNSW IoT attack flows & Consider the data as released with the SOSR 2019 paper: \path{*/\*_flowstats.csv} files representing 10 distinct devices, along with inputs for associated annotations stored under \path{flowdata/flowdata/} - processed to fit our \gref{gs::G1}/\gref{gs::G2} feature schemas. See \texttt{unsw.py}. & \href{https: //iotanalytics.unsw.edu.au/attack-data.html}{UNSW IoT analytics project} provides per-minute "MUD" (metadata about device) flow information + labels for various attacks. No licence found explicitly stated.\\ 
\end{longtable}} 

We make no attempt to verify download date/time/version (except as noted) or any hash of the raw data sources for the latter two datasets (nor for all raw data), and have no further provenance information for the versions of any of the archive releases. This does not impact the provenance information on our frozen campaign.

\Needspace{6\baselineskip} 
\paragraph{Data pre-processing via adapter and filtering/sampling} 
All adapter implementations incorporate a log sign transform of input features, prior to normalisation across the whole dataset pool. Filtering/subsampling procedures detailed below apply to the *remaining* pool of data. For IoT-23, we subsample to a maximum of 2e5 benign / 6e4 attack samples per original individual data bucket file. Furthermore, if the bucket contains >20K rows, we uniformly reduce to 2\% of rows in 1M steps (if the total original file has >1.5G). TON\_IoT is similarly limited to 1.5e5 benign samples (per host) and 6e4 attack samples (per host \& attack-type); furthermore, we similarly uniformly subsample to 3\% if >20K rows (per host \& attack-type) in 2M steps (with a max. sample of 6e4 retained, i.e.\ effectively applying a reservoir subsample first). Reduction counts are rounded-down. CTU-13 is restricted to $\leq$2e5 "normal" rows and 6e4 attack rows per scenario. Gotham similarly limits to $\leq$2e5 benign rows and 4e4 rows per attack-type. N-BaIoT undergoes exact line de-duplication, while respecting the data ordering as timestamps. The UNSW attack flows data represents extracted flows from the dataset at one-minute intervals along with annotations; we do not subsample further.